\unit{Exposição 190 -- Técnica de descida I}
\sid{190}
\src{299}
\seminar{12.\textsuperscript{o} ano, 1959/60, n.\textsuperscript{o} 190\\Dezembro de 1959}
\paperhead{TÉCNICA DE DESCIDA E TEOREMAS DE EXISTÊNCIA EM GEOMETRIA ALGÉBRICA\\
I. GENERALIDADES. DESCIDA POR MORFISMOS FIELMENTE PLANOS}%
  {por Alexander GROTHENDIECK}

Do ponto de vista técnico, a presente exposição e as que se seguirão podem
ser consideradas variações sobre o célebre «teorema 90» de Hilbert.
Parece que A. Weil introduziu o método de descida em geometria algébrica sob
o nome de «descida do corpo base». Aliás, Weil restringiu-se ao caso de
extensões finitas separáveis de corpos. P. Cartier tratou depois o caso de
extensões radiciais de altura 1. Por falta da linguagem dos esquemas e,
mais concretamente, por não dispor de elementos nilpotentes nos anéis que
lhe era permitido considerar, Cartier não pudera formular a identidade
essencial entre esses dois casos.

Atualmente, parece que a técnica geral de descida que será exposta (unida,
quando couber, aos teoremas fundamentais da «geometria formal», cf. [3])
está na base da maioria dos teoremas de existência em geometria
algébrica.\footnote{A errata de 1962 considera excessiva essa formulação: ela é
válida em grande medida para as técnicas não projetivas das duas primeiras
exposições da série TDTE I a VI, mas não para as técnicas projetivas (TDTE IV,
V e VI).} Convém assinalar, além disso, que essa técnica pode, sem dúvida, ser
transposta à «geometria analítica», e cabe esperar que, num futuro próximo,
os especialistas saibam demonstrar os análogos «analíticos» dos teoremas de
existência em geometria formal que serão apresentados na exposição II. Em
todo caso, os trabalhos recentes de Kodaira--Spencer, cujos métodos parecem
inadequados para definir e estudar «variedades de módulos» numa vizinhança
de seus pontos singulares, mostram com bastante clareza a necessidade de
métodos mais próximos da teoria dos esquemas (que deverão completar,
naturalmente, as técnicas transcendentes).

Na presente exposição I trataremos o caso de descida mais elementar (indicado
no título). No entanto, as aplicações dos teoremas 1, 2 e 3 que se seguem já
são muito numerosas. Limitamo-nos a apresentar algumas a título de exemplo,
sem procurar conferir-lhes aqui a máxima generalidade possível.

Utilizaremos livremente a linguagem dos esquemas, para a qual remetemos à
exposição já citada, bem como a [2]. Assinalemos, contudo, expressamente que os
pré-esquemas considerados na presente exposição não são necessariamente
noetherianos, e os

\src{300}
morfismos não são necessariamente de tipo finito. Assim, se $A$ é um anel local
noetheriano, com completamento $\widehat A$, será preciso considerar o anel não
noetheriano $\widehat A\otimes_A\widehat A$, e os morfismos de esquemas afins
correspondentes às inclusões entre os anéis considerados.

\fsection{A}{Preliminares sobre as categorias}

\subsection*{1. Categorias fibradas, dados de descida, morfismos de
$\mathcal F$-descida}
\addcontentsline{toc}{subsection}{1. Categorias fibradas e dados de descida}

\result{a. DEFINIÇÃO}{1.1} Uma \emph{categoria fibrada} $\mathcal F$ de base
$\mathcal C$ consta de uma categoria $\mathcal C$; para cada
$X\in\mathcal C$, de uma categoria $\mathcal F_X$; para cada
$\mathcal C$-morfismo $f:X\longrightarrow Y$, de um funtor
$f^*:\mathcal F_Y\longrightarrow\mathcal F_X$, denotado também por
\[
  f^*(\xi)=\xi\times_Y X
\]
para $\xi\in\mathcal F_Y$ (considerando $X$ como um «objeto de
$\mathcal C$ sobre $Y$», isto é, provido do morfismo $f$); e, por fim,
para dois morfismos componíveis
$X\xrightarrow{f}Y\xrightarrow{g}Z$, de um isomorfismo de funtores
\[
  c_{f,g}:(gf)^*\longrightarrow f^*g^*,
\]
estando esses dados sujeitos às condições seguintes:

\begin{enumerate}
\item[(i)] $\operatorname{id}_X^*=\operatorname{id}$;
\item[(ii)] $c_{f,g}$ é o isomorfismo identidade se $f$ ou $g$ é um
isomorfismo identidade;
\item[(iii)] dados três morfismos componíveis
$X\xrightarrow{f}Y\xrightarrow{g}Z\xrightarrow{h}T$, o seguinte diagrama de
isomorfismos, formado mediante isomorfismos da forma $c_{u,v}$, é
comutativo:
\[
\begin{array}{ccc}
  (h(gf))^*=((hg)f)^* & \longrightarrow & f^*(hg)^* \\
  \big\downarrow && \big\downarrow \\
  (gf)^*h^* & \longrightarrow & (f^*g^*)h^*=f^*(g^*h^*).
\end{array}
\]
\end{enumerate}

\result{EXEMPLO}{1} Seja $\mathcal C$ uma categoria na qual existem os
produtos fibrados. Define-se então uma categoria fibrada $\mathcal F$ de
base $\mathcal C$ tomando $\mathcal F_X=$ categoria dos objetos de
$\mathcal C$ sobre $X$; o funtor
$f^*:\mathcal F_Y\longrightarrow\mathcal F_X$ correspondente a um morfismo
$f:X\longrightarrow Y$ define-se por meio do produto fibrado
$Z\longmapsto Z\times_YX$.

\result{EXEMPLO}{2} Seja $\mathcal C$ a categoria dos pré-esquemas e, para
$X\in\mathcal C$, seja $\mathcal F_X$ a categoria dos feixes
quase-coerentes de módulos sobre $X$. Se $f:X\longrightarrow Y$ é um morfismo
de pré-esquemas, $f^*:\mathcal F_Y\longrightarrow\mathcal F_X$ é o funtor
imagem inversa

\src{301}
de feixes de módulos. Obtém-se assim uma categoria fibrada de base
$\mathcal C$.

\result{b. DEFINIÇÃO}{1.2} Um diagrama
\[
  E\xrightarrow{u}E'
  \overset{v_1}{\underset{v_2}{\rightrightarrows}}E''
\]
de aplicações de conjuntos chama-se \emph{exato} se $u$ é uma bijeção de
$E$ sobre a parte de $E'$ formada pelos $x'\in E'$ tais que
$v_1(x')=v_2(x')$.

\result{DEFINIÇÃO}{1.3} Seja $\mathcal F$ uma categoria fibrada de base
$\mathcal C$; consideremos um diagrama
\[
  \tag{+}
  S\xleftarrow{\alpha}S'
  \overset{\beta_1}{\underset{\beta_2}{\rightrightarrows}}S''
\]
de morfismos em $\mathcal C$ tais que
$\alpha\beta_1=\alpha\beta_2$; esse diagrama chama-se
$\mathcal F$-\emph{exato} se, para cada par $\xi,\eta$ de elementos de
$\mathcal F_S$, o seguinte diagrama de aplicações de conjuntos
\[
  \tag{+'}
  \operatorname{Hom}(\xi,\eta)
  \xrightarrow{\alpha^*}
  \operatorname{Hom}(\alpha^*\xi,\alpha^*\eta)
  \overset{\beta_1^*}{\underset{\beta_2^*}{\rightrightarrows}}
  \operatorname{Hom}(\gamma^*\xi,\gamma^*\eta),
\]
onde $\gamma=\alpha\beta_1=\alpha\beta_2$, é exato. Neste último diagrama
temos identificado $\beta_i^*\alpha^*$ com
$(\alpha\beta_i)^*=\gamma^*$ mediante $c_{\beta_i,\alpha}$, para simplificar.

\result{DEFINIÇÃO}{1.4} Seja $\mathcal F$ uma categoria fibrada de base
$\mathcal C$; consideremos dois morfismos $\beta_1,\beta_2$ de $S''$ em $S'$,
em $\mathcal C$. Seja
$\xi'\in\mathcal F_{S'}$. Chama-se \emph{dado de colagem} sobre $\xi'$
(relativo ao par $(\beta_1,\beta_2)$) a um isomorfismo de
$\beta_1^*(\xi')$ sobre $\beta_2^*(\xi')$. Se $\xi',\eta'\in\mathcal F_{S'}$
estão munidos, cada um, de um dado de colagem, um morfismo
$u:\xi'\longrightarrow\eta'$ em $\mathcal F_{S'}$ se diz
\emph{compatível com os dados de colagem} se o diagrama seguinte é
comutativo:
\[
\begin{array}{ccc}
  \beta_1^*(\xi') & \longrightarrow & \beta_2^*(\xi') \\
  \big\downarrow && \big\downarrow \\
  \beta_1^*(\eta') & \longrightarrow & \beta_2^*(\eta').
\end{array}
\]
Desse modo, os objetos de $\mathcal F_{S'}$ providos de um dado de colagem
(relativo a $\beta_1,\beta_2$) formam uma categoria. Se
$\alpha:S'\longrightarrow S$ é um morfismo tal que
$\alpha\beta_1=\alpha\beta_2$, então, para todo $\xi\in\mathcal F_S$, o
objeto $\xi'=\alpha^*(\xi)$

\src{302}
de $\mathcal F_{S'}$ está provido de um dado de colagem canônico, visto que
\[
  \beta_i^*\alpha^*(\xi)\simeq(\alpha\beta_i)^*(\xi)=\gamma^*(\xi),
\]
tomando de novo $\gamma=\alpha\beta_1=\alpha\beta_2$; além disso, se
$u:\xi\longrightarrow\eta$ é um morfismo em $\mathcal F_S$, então
\[
  \alpha^*(u):\alpha^*(\xi)\longrightarrow\alpha^*(\eta)
\]
é um morfismo em $\mathcal F_{S'}$ compatível com os dados de colagem
canônicos. Obtém-se assim um funtor canônico da categoria $\mathcal F_S$ para
a categoria dos objetos de $\mathcal F_{S'}$ providos de um dado de colagem
relativo ao par $(\beta_1,\beta_2)$. Dito isto, a definição 1.3 também
pode ser expressa dizendo que o diagrama (+) dessa definição é
$\mathcal F$-exato se o funtor anterior é «plenamente fiel», isto é, se
define uma equivalência da categoria $\mathcal F_S$ com uma subcategoria da
categoria dos objetos de $\mathcal F_{S'}$ providos de um dado de colagem
relativo a $(\beta_1,\beta_2)$.

\result{DEFINIÇÃO}{1.5} Diz-se que um dado de colagem sobre
$\xi'\in\mathcal F_{S'}$ é \emph{efetivo} (com respeito a $\alpha$) se
$\xi'$, provido deste dado, é isomorfo a um $\alpha^*(\xi)$, com
$\xi\in\mathcal F_S$.

Quando o diagrama (+) é $\mathcal F$-exato, o objeto $\xi$ da definição
anterior fica, portanto, determinado salvo isomorfismo único, e a categoria
$\mathcal F_S$ é equivalente à categoria dos objetos de
$\mathcal F_{S'}$ providos de um dado de colagem efetivo.

\medskip\noindent
c. O caso mais importante é aquele em que
\[
  S''=S'\times_S S',
\]
sendo as $\beta_i$ as duas projeções $p_1$ e $p_2$ de
$S'\times_S S'$ sobre seus dois fatores (supomos agora que em
$\mathcal C$ existem os produtos fibrados). Têm-se então duas
condições necessárias imediatas para que seja efetivo um dado de colagem
$\varphi:p_1^*(\xi')\longrightarrow p_2^*(\xi')$ sobre um
$\xi'\in\mathcal F_{S'}$. Tem-se
\[
  \tag{i}
  \Delta^*(\varphi)=\operatorname{id}_{\xi'},
\]
onde $\Delta:S'\longrightarrow S'\times_S S'$ designa o morfismo diagonal,
e onde se identificou $\Delta^*p_i^*(\xi')$ com
$(p_i\Delta)^*(\xi')=\xi'$. Tem-se
\[
  \tag{ii}
  p_{31}^*(\varphi)=p_{32}^*(\varphi)\,p_{21}^*(\varphi),
\]

\src{303}
onde $p_{ij}$ designa a projeção canônica de
$S'\times_S S'\times_S S'$ sobre o produto parcial de seus fatores
$i$-ésimo e $j$-ésimo.

\result{DEFINIÇÃO}{1.6} Chama-se \emph{dado de descida} sobre
$\xi'\in\mathcal F_{S'}$, relativo ao morfismo
$\alpha:S'\longrightarrow S$, a um dado de colagem sobre $\xi'$ relativo ao par
$(p_1,p_2)$ de projeções canônicas
$S'\times_S S'\longrightarrow S'$ que satisfaz as condições (i) e (ii)
anteriores.

\result{DEFINIÇÃO}{1.7} Um morfismo $\alpha:S'\longrightarrow S$ chama-se
\emph{morfismo de $\mathcal F$-descida} se o diagrama de morfismos
\[
  S\xleftarrow{\alpha}S'
  \overset{p_1}{\underset{p_2}{\rightrightarrows}}S'\times_S S'
\]
é $\mathcal F$-exato (definição 1.3); diz-se que $\alpha$ é um
\emph{morfismo de $\mathcal F$-descida estrito} se, além disso, todo dado de
descida (definição 1.6) sobre um objeto de $\mathcal F_{S'}$ é efetivo.

Esta última condição também pode ser enunciada, de modo mais gráfico, assim:
«Dá no mesmo» dar um objeto de $\mathcal F_S$ ou um objeto de
$\mathcal F_{S'}$ provido de um dado de descida.

Observe-se que se um morfismo $\alpha:S'\longrightarrow S$ admite uma «seção»
$s:S\longrightarrow S'$ (isto é, um morfismo $s$ tal que
$\alpha s=\operatorname{id}_S$), então é um morfismo de
$\mathcal F$-descida estrito.\footnote{A errata de 1962 esclarece que não é
necessário supor previamente que $\alpha$ seja um morfismo de
$\mathcal F$-descida.} Com efeito, se $\xi'\in\mathcal F_{S'}$ está provido de
um dado de descida relativo a $\alpha$, «provém» de
$\xi=s^*(\xi')$.

\medskip\noindent
d. As noções anteriores podem ser apresentadas de maneira mais intuitiva
introduzindo, para um objeto parâmetro $T$ de $\mathcal C$ sobre $S$, o
conjunto
\[
  \operatorname{Hom}_S(T,S')=S'(T),
\]
cujos elementos serão denotados por $t,t'$, etc. Dado um objeto
$\xi'\in\mathcal F_{S'}$, a cada $t\in S'(T)$ corresponde então um objeto
$t^*(\xi')$ de $\mathcal F_T$, que também será denotado por $\xi'_t$. Um dado de
colagem sobre $\xi'$ (relativo a $p_1,p_2$) fica definido então dando, para
todo $T$ sobre $S$ e todo par de pontos $t,t'\in S'(T)$, um isomorfismo
\[
  \varphi_{t',t}:\xi'_t\longrightarrow\xi'_{t'}
\]
(que satisfaz condições evidentes de funtorialidade ao variar $T$). As
condições (i) e (ii) de 1.c escrevem-se então
\[
  \tag{i bis}
  \varphi_{t,t}=\operatorname{id}_{\xi'_t}
  \qquad\text{para todo $T$ e todo $t\in S'(T)$},
\]

\src{304}
\[
  \tag{ii bis}
  \varphi_{t,t''}=\varphi_{t,t'}\varphi_{t',t''}
  \qquad\text{para todo $T$ e todos $t,t',t''\in S'(T)$}.
\]
Vê-se, além disso, que (ii bis), tomando $t=t'=t''$, implica
$\varphi_{t,t}=\varphi_{t,t}^2$; visto que $\varphi_{t,t}$ é um isomorfismo
por hipótese, daí se obtém a relação (i bis), que é de fato uma consequência
de (ii bis) (portanto, (i) é consequência de (ii)). Mas, se já não se supõe
a priori que os $\varphi_{t,t'}$ sejam isomorfismos (isto é, que
$\varphi:p_1^*(\xi')\longrightarrow p_2^*(\xi')$ seja um isomorfismo), então
(ii bis) não implica necessariamente (i bis); no entanto, a conjunção de
(ii bis) e (i bis) implica que os $\varphi_{t,t'}$ são isomorfismos (pois se
terá
$\varphi_{t,t'}\varphi_{t',t}=\varphi_{t,t}=\operatorname{id}_{\xi'_t}$).

\subsection*{2. Diagramas exatos e epimorfismos estritos, morfismos de
descida. Exemplos}
\addcontentsline{toc}{subsection}{2. Diagramas exatos e epimorfismos estritos}

\result{a. DEFINIÇÃO}{2.1} Seja $\mathcal C$ uma categoria. Um diagrama
\[
  T\xrightarrow{\alpha}T'
  \overset{\beta_1}{\underset{\beta_2}{\rightrightarrows}}T''
\]
de morfismos em $\mathcal C$ chama-se \emph{exato} se, para todo
$Z\in\mathcal C$, o diagrama correspondente de aplicações de conjuntos
\[
  \operatorname{Hom}(Z,T)\longrightarrow
  \operatorname{Hom}(Z,T')
  \rightrightarrows\operatorname{Hom}(Z,T'')
\]
é exato (definição 1.2). Diz-se então que $(T,\alpha)$ (ou, por abuso de
linguagem, que $T$) é um \emph{núcleo} do par de morfismos
$(\beta_1,\beta_2)$.

Evidentemente, esse núcleo está determinado salvo isomorfismo único. Se
$\mathcal C$ é a categoria dos conjuntos, a definição anterior é
compatível com a definição 1.2. Define-se de maneira dual a exatidão de um
diagrama
\[
  S\xleftarrow{\alpha}S'
  \overset{\beta_1}{\underset{\beta_2}{\rightrightarrows}}S''
\]
de morfismos em $\mathcal C$; diz-se então que $(S,\alpha)$ é um
\emph{conúcleo} do par de morfismos $(\beta_1,\beta_2)$.

\result{DEFINIÇÃO}{2.2} Um morfismo $\alpha:S'\longrightarrow S$ chama-se
\emph{epimorfismo estrito} se é um epimorfismo e se, para todo morfismo
$u:S'\longrightarrow Z$, a seguinte condição necessária é também
suficiente para que $u$ se fatore como
$S'\longrightarrow S\longrightarrow Z$: para todo $S''\in\mathcal C$ e todo
par de morfismos $\beta_1,\beta_2:S''\longrightarrow S'$ tal que
$\alpha\beta_1=\alpha\beta_2$, tem-se também $u\beta_1=u\beta_2$.

Se existe o produto fibrado $S'\times_S S'$, equivale a dizer que o diagrama
\[
  S\xleftarrow{\alpha}S'
  \overset{p_1}{\underset{p_2}{\rightrightarrows}}S'\times_S S'
\]

\src{305}
é exato, isto é, que $S$ é um conúcleo do par $(p_1,p_2)$. Em todo
caso, um morfismo conúcleo é um epimorfismo estrito. Assinalemos também que um
epimorfismo estrito que seja um monomorfismo é um isomorfismo. Deixamos ao
leitor desenvolver a noção dual de \emph{monomorfismo estrito}.

Para precisar as relações entre a noção de morfismo de
$\mathcal F$-descida e a de epimorfismo estrito, introduzimos além disso as
definições seguintes:

\result{DEFINIÇÃO}{2.3} Um morfismo $\alpha:S'\longrightarrow S$ chama-se
\emph{epimorfismo universal} (resp. \emph{epimorfismo estrito universal}) se,
para todo $T$ sobre $S$, existe o produto fibrado $T'=S'\times_ST$ e a
projeção $T'\longrightarrow T$ é um epimorfismo (resp. um epimorfismo
estrito).

Nas categorias muito boas (como a categoria dos conjuntos, a
categoria dos feixes de conjuntos sobre um espaço topológico, as categorias
abelianas, etc.), as quatro noções de epimorfismo assim introduzidas
coincidem; ao contrário, todas elas são distintas numa categoria como
a dos pré-esquemas, ou a dos pré-esquemas sobre um pré-esquema não vazio
dado $S$, inclusive se nos limitamos a $S$-esquemas finitos sobre $S$.

\result{DEFINIÇÃO}{2.4} Um morfismo $\alpha:S'\longrightarrow S$ chama-se
\emph{morfismo de descida} (resp. \emph{morfismo de descida estrito}) se é
um morfismo de $\mathcal F$-descida (resp. um morfismo de
$\mathcal F$-descida estrito) (cf. definição 1.7), onde $\mathcal F$ denota
a categoria fibrada de base $\mathcal C$ dos objetos de $\mathcal C$ sobre
objetos de $\mathcal C$ (n\textsuperscript{o} 1, exemplo 1).

\result{PROPOSIÇÃO}{2.1} Se em $\mathcal C$ existem os produtos finitos e
os produtos fibrados (finitos), então, em $\mathcal C$, os morfismos de
descida coincidem com os epimorfismos estritos universais.

\result{b. EXEMPLOS}{} Seja $\mathcal C$ a categoria dos pré-esquemas. Seja
$S\in\mathcal C$, e sejam $S'$ e $S''$ dois pré-esquemas finitos sobre $S$,
isto é, correspondentes a feixes de álgebras $\mathcal A'$, $\mathcal A''$
sobre $S$ que, como feixes de módulos, são quase-coerentes e de tipo finito
(isto é, coerentes se $S$ é localmente noetheriano). Seja
$\alpha:S'\longrightarrow S$ o morfismo estrutural de $S'$, e sejam
$\beta_1,\beta_2$ dois $S$-morfismos de $S''$ em $S'$, definidos por
homomorfismos de álgebras $\mathcal A'\longrightarrow\mathcal A''$, denotados
também por $\beta_1,\beta_2$. Utilizando o fato de que um morfismo finito é
fechado (primeiro teorema de Cohen--Seidenberg), demonstra-se facilmente que o
diagrama
\[
  \tag{+}
  S\xleftarrow{\alpha}S'
  \overset{\beta_1}{\underset{\beta_2}{\rightrightarrows}}S''
\]
em $\mathcal C$ é exato se e somente se também o é o diagrama de feixes

\src{306}
\[
  \mathcal O_S=\mathcal A\longrightarrow\mathcal A'
  \overset{\beta_1}{\underset{\beta_2}{\rightrightarrows}}\mathcal A''
\]
sobre $S$. Em particular, se $\alpha:S'\longrightarrow S$ é um morfismo
finito correspondente a um feixe $\mathcal A'$ de álgebras sobre $S$, então
$\alpha$ é um epimorfismo estrito se e somente se o diagrama de feixes
\[
  \mathcal O_S=\mathcal A\longrightarrow\mathcal A'
  \overset{p_1}{\underset{p_2}{\rightrightarrows}}
  \mathcal A'\otimes_{\mathcal A}\mathcal A'
\]
é exato (é um epimorfismo se e somente se
$\mathcal A\longrightarrow\mathcal A'$ é injetivo). Se $S$ é afim de anel
$A$, e portanto $S'$ é afim de anel $A'$ finito sobre $A$, então
$S'\longrightarrow S$ é um epimorfismo estrito se e somente se
$A\longrightarrow A'$ é um isomorfismo de $A$ sobre o subanel de $A'$
formado pelos $x'\in A'$ tais que
\[
  1_{A'}\otimes_Ax'-x'\otimes_A1_{A'}=0
\]
(é um epimorfismo se e somente se $A\longrightarrow A'$ é injetivo). Como já
assinalamos, mesmo se $S$ é o esquema de um anel local artiniano, um
morfismo finito $S'\longrightarrow S$ que seja um epimorfismo não é
necessariamente um epimorfismo estrito. No entanto, pode-se demonstrar que, se
$S$ é um pré-esquema noetheriano, todo morfismo finito
$S'\longrightarrow S$ que seja um epimorfismo é a composição de uma sucessão
finita de epimorfismos estritos (também finitos). Isso mostra, além disso, que
a composição de dois epimorfismos estritos não é necessariamente um
epimorfismo estrito.

\result{c}{} Se $(+)$ é um diagrama exato de morfismos finitos, então,
para todo morfismo plano $T\longrightarrow S$ de pré-esquemas, o diagrama
obtido de $(+)$ mediante a mudança de base $T\longrightarrow S$ continua
exato. Disso se deduz que, se $X,Y$ são dois $S$-pré-esquemas, sendo $X$
plano sobre $S$, então o seguinte diagrama de aplicações de conjuntos
(onde $X',Y'$ são as imagens inversas de $X,Y$ sobre $S'$, e $X'',Y''$ suas
imagens inversas sobre $S''$)
\[
  \operatorname{Hom}_S(X,Y)\longrightarrow
  \operatorname{Hom}_{S'}(X',Y')
  \rightrightarrows
  \operatorname{Hom}_{S''}(X'',Y'')
\]
é exato. Em particular, se $\mathcal F$ denota a categoria fibrada cuja base
é a categoria $\mathcal C$ dos pré-esquemas e tal que, para
$X\in\mathcal C$, $\mathcal F_X$ é a categoria dos $X$-pré-esquemas
planos, então o diagrama $(+)$ é $\mathcal F$-exato. (Esse resultado deixa
de ser verdadeiro se não se impõe a hipótese de planicidade; em particular, um
epimorfismo estrito finito não é necessariamente um morfismo de descida.)
Vê-se igualmente que $(+)$ é $\mathcal F$-exato se $\mathcal F$ denota a
categoria fibrada para a qual $\mathcal F_X$ é a categoria dos feixes
quase-coerentes e planos sobre o pré-esquema $X$ (também aqui a hipótese

\src{307}
de planicidade é essencial). Num e noutro caso, a questão da efetividade de
um dado de colagem (e, mais concretamente, de um dado de descida quando
$S''=S'\times_SS'$) sobre um objeto plano sobre $S'$ é delicada (e sua
resposta em diversos casos particulares é um dos objetos principais das
presentes exposições). O conferenciante ignora se, para todo epimorfismo
estrito finito $S'\longrightarrow S$, todo dado de descida sobre um feixe
quase-coerente plano sobre $S'$ é efetivo (mesmo supondo que $S$ seja o
espectro de um anel local artiniano e limitando-nos aos feixes localmente
livres de posto 1). Mais geralmente, sejam $A$ um anel e $A'$ uma
$A$-álgebra (tudo é comutativo) tal que o diagrama de aplicações
\[
  A\longrightarrow A'\rightrightarrows A'\otimes_AA'
\]
seja exato, o que equivale também a que o morfismo
$S'\longrightarrow S$ correspondente entre os espectros de $A,A'$ seja um
morfismo de $\mathcal F$-descida, onde $\mathcal F$ é a categoria fibrada
dos feixes quase-coerentes planos. Seja $M'$ um $A'$-módulo plano provido de
um dado de descida a $A$, isto é, de um isomorfismo
\[
  \varphi:M'\otimes_AA'\xrightarrow{\sim}A'\otimes_AM'
\]
de $A'\otimes_AA'$-módulos que satisfaz as condições (i) e (ii) de 1.c
(cuja explicitação em termos de módulos deixamos ao leitor). É efetivo
este dado (com respeito à categoria fibrada dos feixes quase-coerentes
planos)? Seja $M$ o subconjunto de $M'$ formado pelos $x'\in M'$ tais que
\[
  \varphi(x'\otimes_A1_{A'})=1_{A'}\otimes_Ax',
\]
que é um sub-$A$-módulo de $M'$. A injeção canônica
$M\longrightarrow M'$ define um homomorfismo de $A'$-módulos
$M\otimes_AA'\longrightarrow M'$. A efetividade de $\varphi$ significa
então o seguinte: $M$ é um $A$-módulo plano e o homomorfismo anterior é
um isomorfismo.

\result{OBSERVAÇÃO}{} Nas considerações anteriores não havíamos imposto
nenhuma hipótese de planicidade aos morfismos do diagrama $(+)$, o que nos
obrigava, para dispor de uma técnica de descida, a impor hipóteses de
planicidade aos objetos sobre $S,S'$ considerados. No parágrafo 2 imporemos
uma hipótese de planicidade a $\alpha:S'\longrightarrow S$, o que nos permitirá
dispor de uma técnica de descida para objetos sobre $S,S'$ que já não estão
submetidos a condição alguma de planicidade. Em todos os casos intervém uma
hipótese de planicidade. Essa é uma das principais razões da importância da
noção de planicidade em geometria algébrica (cujo papel não podia manifestar-se
enquanto nos limitávamos a corpos base, sobre os quais, com efeito, tudo é
plano!).

\src{308}
\subsection*{3. Aplicação aos morfismos étale}
\addcontentsline{toc}{subsection}{3. Aplicação aos morfismos étale}

Sejam $A$ um anel local e $B$ um álgebra local sobre $A$ cujo ideal maximal
induz o de $A$. Diremos que $B$ é \emph{étale} sobre $A$ (em lugar de
«não ramificada», expressão usada em outros lugares) se satisfaz as condições
seguintes:
\begin{enumerate}
\item[(i)] $B$ é plana sobre $A$;
\item[(ii)] $B/\mathfrak mB$ é uma extensão finita separável de
$A/\mathfrak m=k$ (onde $\mathfrak m$ denota o ideal maximal de $A$).
\end{enumerate}
Quando $A$ e $B$ são noetherianos e $k$ é algebricamente fechado, isso
significa que o homomorfismo $\widehat A\longrightarrow\widehat B$ entre os
completamentos que prolonga $A\longrightarrow B$ é um isomorfismo. Um morfismo de
tipo finito $f:T\longrightarrow S$ diz-se \emph{étale em $x\in T$}, ou também
que $T$ é \emph{étale sobre $S$ em $x$}, se $\mathcal O_x$ é étale sobre
$\mathcal O_{f(x)}$; diz-se que $f$ é \emph{étale}, ou chama-se $f$ um
\emph{morfismo étale}, ou diz-se que $T$ é \emph{étale sobre $S$}, se $f$ é
étale em todo $x\in T$. Assinalemos, além disso, que se $S$ é localmente
noetheriano, o conjunto de pontos de $T$ nos quais $f$ é étale é aberto, e
o uso do «main theorem» de Zariski permite precisar a estrutura de $T/S$
numa vizinhança de tal ponto (por meio de uma equação de tipo bem conhecido).

Quando $S$ é um esquema de tipo finito sobre o corpo dos números
complexos, corresponde-lhe um espaço analítico $\overline S$ no sentido de
Serre [5], exceto que $\overline S$ pode conter elementos nilpotentes em seu
feixe estrutural, o que em nada altera o essencial de [5]. Vê-se então facilmente
que $f$ é étale se e somente se
$\overline f:\overline T\longrightarrow\overline S$ também o é, isto é, se todo
ponto de $\overline T$ admite uma vizinhança na qual $\overline f$ induz um
isomorfismo sobre um aberto de $\overline S$. Em particular, a todo
\emph{recobrimento étale} $T$ de $S$ (isto é, um morfismo finito étale
$f:T\longrightarrow S$) corresponde um recobrimento étale $\overline T$ de
$\overline S$, que é conexo se e somente se $T$ também o é [5]. Vê-se, além
disso, facilmente que, se $T,T'$ são dois esquemas étale sobre $S$, então a
aplicação natural
\[
  \operatorname{Hom}_S(T,T')\longrightarrow
  \operatorname{Hom}_{\overline S}(\overline T,\overline{T'})
\]
é bijetiva, isto é, o funtor $T\longmapsto\overline T$ da categoria dos
esquemas étale sobre $S$ à categoria dos espaços analíticos étale
sobre $\overline S$ é «plenamente fiel» e, portanto, define uma equivalência
da primeira categoria com uma subcategoria da segunda. Um teorema de
Grauert--Remmert [2] implica que, se $S$ é normal, obtém-se assim uma
equivalência entre a categoria dos \emph{recobrimentos étale} de $S$ e a
categoria dos recobrimentos étale (finitos) de $\overline S$, isto é, que
todo recobrimento étale de $\overline S$ é $\overline S$-isomorfo a um
$\overline T$, onde $T$ é um recobrimento étale de $S$. Demonstremos que o
\emph{teorema de Grauert--Remmert continua válido sem hipótese de
normalidade sobre $S$}. Seja, com efeito, primeiro $S'\longrightarrow S$ um
epimorfismo estrito finito; suponhamos demonstrado o

\src{309}
teorema para $S'$ e demonstremos que será verdadeiro para $S$. Seja
$\mathfrak T$ um recobrimento étale de $\overline S$ e consideremos sua imagem
inversa $\mathfrak T'$ sobre $\overline{S'}$, que corresponde a um feixe
analítico coerente $\mathcal A'$ de álgebras sobre $\overline{S'}$, imagem
inversa do feixe de álgebras $\mathcal A$ sobre $\overline S$ que define
$\mathfrak T$. Por hipótese, sobre $\overline{S'}$, $\mathfrak T'$ provém de
um recobrimento étale $T'$ de $S'$, isto é, $\mathcal A'$ provém de um feixe
coerente de álgebras $A'$ sobre $S'$. Por outro lado, $\mathcal A'$ está
provido de um dado de descida canônico relativo a
$\overline{S'}\longrightarrow\overline S$, isto é, de um isomorfismo entre
suas duas imagens inversas sobre
\[
  \overline{S'}\times_{\overline S}\overline{S'}
  =\overline{(S'\times_SS')},
\]
(que satisfaz as condições (i), (ii)), e esse isomorfismo provém, conforme
dito, de um isomorfismo entre os feixes algébricos correspondentes, isto é,
de um dado de descida sobre $A'$ relativo a
$S'\longrightarrow S$. Verifica-se facilmente que este último é efetivo
(pois o de $\mathcal A'$ também o é, e a efetividade de um dado de descida,
tal como se explicitou no número anterior, reconhece-se localmente nos
\emph{completamentos} dos módulos que intervêm). Obtém-se assim um feixe
coerente de álgebras $A$ sobre $S$, que define um recobrimento $T$ de $S$,
o recobrimento procurado. O resultado anterior continua então manifestamente
válido se $S'\longrightarrow S$ for apenas uma composição de um número finito
de epimorfismos estritos finitos, isto é, se for um epimorfismo finito
qualquer (pelo resultado de fatoração assinalado no parágrafo 2). Segue-se que
o teorema de Grauert--Remmert continua válido se $S$ é um esquema
\emph{reduzido}, isto é, se $\mathcal O_S$ não tem elementos nilpotentes, como
se vê introduzindo sua normalização $S'$. Daí se passa facilmente ao caso geral.

Uma demonstração inteiramente análoga, que utiliza de novo o resultado de
fatoração para os epimorfismos estritos finitos e a natureza «formal» da
efetividade dos dados de descida, permite demonstrar o resultado seguinte:
seja $S$ um pré-esquema localmente noetheriano e seja
$S'\longrightarrow S$ um morfismo \emph{finito, sobrejetivo e radicial} (ou,
o que dá no mesmo, um morfismo de tipo finito tal que, para todo $T$
sobre $S$, o morfismo $T'=S'\times_ST\longrightarrow T$ seja um homeomorfismo,
o que também se exprime dizendo que $S'\longrightarrow S$ é um
\emph{homeomorfismo universal}). Para todo $T$ étale sobre $S$, consideremos
sua imagem inversa $T'=T\times_SS'$, que é étale sobre $S'$. Então o
funtor $T\longmapsto T'$ é uma \emph{equivalência de categorias entre a
categoria dos pré-esquemas $T$ étale sobre $S$ e a categoria dos
pré-esquemas $T'$ étale sobre $S'$}. (Utiliza-se a bijetividade de
\[
  \operatorname{Hom}_S(T_1,T_2)\longrightarrow
  \operatorname{Hom}_{S'}(T'_1,T'_2)
\]
para dois pré-esquemas $T_1,T_2$ étale sobre $S$, fato cuja verificação direta
é fácil, e o fato de que o teorema enunciado é verdadeiro se
$S'=(S,\mathcal O_S/\mathcal J)$,

\src{310}
onde $\mathcal J$ é um feixe coerente nilpotente de ideais de $\mathcal O_S$
([4], lema 6)). Observe-se, além disso, que não supomos aqui que os $T$
considerados sejam finitos sobre $S$. Este resultado implica, em particular,
que o morfismo $S'\longrightarrow S$ induz um isomorfismo do grupo
fundamental $[\,]$ de $S'$ sobre o de $S$ («\emph{invariância topológica do
grupo fundamental de um pré-esquema}»).

\subsection*{4. Relações com a 1-cohomologia}
\addcontentsline{toc}{subsection}{4. Relações com a 1-cohomologia}

\medskip\noindent
a. Seja $\mathcal C$ uma categoria na qual sempre existe o produto de dois
objetos, e seja $T\in\mathcal C$. Para todo conjunto finito
$I\ne\varnothing$, pode-se considerar $T^I$; ao variar $I$, obtém-se assim um
funtor covariante da categoria dos conjuntos finitos não vazios para
$\mathcal C$, isto é, o que se pode chamar um \emph{objeto simplicial} de
$\mathcal C$, denotado por $K_T$. Este último depende covariantemente de $T$;
além disso, se $u,v$ são dois morfismos $T\longrightarrow T'$, os morfismos
correspondentes $K_T\longrightarrow K_{T'}$ são homotópicos. Diremos que $T$
\emph{domina} a $T'$ se $\operatorname{Hom}(T,T')\ne\varnothing$; esta é uma
relação de pré-ordem filtrante crescente em $\mathcal C$. Do anterior se
deduz que, se $T$ domina $T'$, existe uma classe canônica (salvo homotopia)
de homomorfismos de objetos simpliciais $K_T\longrightarrow K_{T'}$; em
particular, se $K_T$ e $K_{T'}$ são tais que cada um domina ao outro, então
$K_T$ e $K_{T'}$ são homotopicamente equivalentes. Seja agora $F$ um funtor
(contravariante, para fixar ideias) de $\mathcal C$ para uma categoria abeliana
$\mathcal C'$; então
\[
  C^*(T,F)=F(K_T)
\]
é um objeto cosimplicial de $\mathcal C'$ e, portanto, define da maneira
habitual um complexo (de cocadeias) em $\mathcal C'$, do qual se pode tomar a
cohomologia
\[
  H^*(T,F)=H^*(C^*(T,F))=H^*(F(K_T))
\]
(poder-se-á escrever uma $\mathcal C$ como subíndice de $H^*$ se houver
possibilidade de confusão). Esse é um funtor cohomológico em $F$, cuja
variância ao variar $T$ se deduz do que foi dito sobre os $K_T$; mais
precisamente, para $F$ fixo e $T$ variável em $\mathcal C$ (pré-ordenada pela
relação de dominação), os $H^*(T,F)$ formam um sistema indutivo de objetos
graduados de $\mathcal C'$;
em particular, se $T$ e $T'$ são tais que cada um domina ao outro, então
$H^*(T,F)$ e $H^*(T',F)$ são canonicamente isomorfos.

Suponhamos que em $\mathcal C$ existem os produtos fibrados. Para
$S\in\mathcal C$ fixo, pode-se aplicar o anterior à categoria
$\mathcal C_S$ dos objetos de $\mathcal C$ sobre $S$; escreveremos
$C^*(T/S,F)$ e $H^*(T/S,F)$ em lugar de $C^*(T,F)$ e $H^*(T,F)$ quando
quisermos precisar que trabalhamos na categoria $\mathcal C_S$; assim,

\src{311}
$C^*(T/S,F)$ é um complexo de cocadeias em $\mathcal C'$ que, em dimensão
$n$, é igual a
\[
  F(T\times_ST\times_S\cdots\times_ST)
\]
(onde o parêntese contém $n+1$ fatores).

Observe-se que, como de costume, pode-se definir $H^0(T/S,F)$ sem supor
abeliana a categoria $\mathcal C'$: é o núcleo (definição 2.1), se existe,
do par de morfismos $F(p_i)$ ($i=1,2$)
\[
  F(T)\rightrightarrows F(T\times_ST)
\]
correspondentes às duas projeções
$p_1,p_2:T\times_ST\rightrightarrows T$. Em particular, ter-se-á um morfismo
natural (chamado de \emph{aumento})
\[
  F(S)\longrightarrow H^0(T/S,F)
\]
que será um isomorfismo nos casos favoráveis (em particular, se
$T\longrightarrow S$ é um epimorfismo estrito e $F$ é «exato pela
esquerda»). Do mesmo modo, quando $F$ toma valores na categoria dos
grupos numa categoria $\mathcal C''$, também se pode definir $H^1(T/S,F)$;
quando $\mathcal C''$ é a categoria dos conjuntos (isto é, $F$ toma
valores na categoria dos grupos ordinários, não necessariamente
comutativos), $H^1(T/S,F)$ é o quociente do subgrupo $Z^1(T/S,F)$ de
$C^1(T/S,F)=F(T\times_ST)$ formado pelos $g$ tais que
\[
  F(p_{31})(g)=F(p_{32})(g)\,F(p_{21})(g)
\]
pelo grupo de operadores $F(T)$, que atua sobre $C^1(T/S,F)$ e, em
particular, sobre o subconjunto $Z^1(T/S,F)$ mediante
\[
  \rho(g')\cdot g=F(p_2)(g')\,g\,F(p_1)(g')^{-1}.
\]

\medskip\noindent
b. Seja, por exemplo, $\mathcal F$ uma categoria fibrada de base $\mathcal C$.
Sejam $\xi,\eta\in\mathcal F_S$ e, para todo $S'$ sobre $S$, seja
\[
  F_{\xi,\eta}(S')=
  \operatorname{Hom}(\xi\times_SS',\eta\times_SS').
\]
Assim, $F_{\xi,\eta}$ é um funtor contravariante de $\mathcal C_S$ para a
categoria dos conjuntos. Dito isto, \emph{afirmar que o morfismo de
aumento}
\[
  F_{\xi,\eta}(S)\longrightarrow H^0(S'/S,F_{\xi,\eta})
\]
\emph{é um isomorfismo para todo par de elementos
$\xi,\eta\in\mathcal F_S$ significa que
$\alpha:S'\longrightarrow S$ é um morfismo de $\mathcal F$-descida}
(definição 1.7).

\medskip\noindent
c. Ponhamos igualmente, para $\xi\in\mathcal F_S$ e todo objeto $S'$ de
$\mathcal C$ sobre $S$,
\[
  G_\xi(S')=\operatorname{Aut}(\xi\times_SS') ;
\]
definiu-se assim um funtor contravariante $G_\xi$ de $\mathcal C_S$ para a
categoria dos

\src{312}
grupos. Dito isto, constata-se que $Z^1(S'/S,G_\xi)$ \emph{se identifica
canonicamente com o conjunto dos dados de descida sobre
$\xi'=\xi\times_SS'$ relativos a $S'\longrightarrow S$} (definição 1.6), e
$H^1(S'/S,G_\xi)$ \emph{se identifica com o conjunto das classes (salvo
isomorfismo) de objetos de $\mathcal F_{S'}$ providos de um dado de descida
relativo a $\alpha:S'\longrightarrow S$ que, como objetos de
$\mathcal F_{S'}$, são isomorfos a $\xi'=\xi\times_SS'$}. Portanto, se
$\alpha:S'\longrightarrow S$ é um morfismo de $\mathcal F$-descida
(cf. (b)), então $H^1(S'/S,G_\xi)$ \emph{contém como subconjunto o
conjunto das classes (salvo isomorfismo) de objetos $\eta$ de
$\mathcal F_S$ tais que $\eta\times_SS'$ seja isomorfo em
$\mathcal F_{S'}$ a $\xi\times_SS'$; e essa inclusão é uma identidade se e
somente se todo dado de descida sobre $\xi'=\xi\times_SS'$ relativo a
$\alpha:S'\longrightarrow S$ é efetivo}. (Esse será o caso, em particular,
se $\alpha:S'\longrightarrow S$ é um morfismo de $\mathcal F$-descida
estrito.)

\result{OBSERVAÇÃO}{} Os complexos de cocadeias do tipo $C^*(T/S,F)$
contêm como casos particulares a maioria dos complexos padrão
conhecidos (cohomologia de Čech, cohomologia de grupos, etc.) e desempenham um
papel importante em geometria algébrica, em particular na «cohomologia de
Weil» dos pré-esquemas.

\medskip\noindent
d. \result{EXEMPLO}{1} Seja $S'$ um objeto sobre $S\in\mathcal C$, e seja
$\Gamma$ um grupo de automorfismos de $S'$ tal que $S'$ seja «formalmente
principal sobre $S$, de grupo $\Gamma$», isto é, tal que o morfismo natural
\[
  \Gamma\times S'\longrightarrow S'\times_SS',
\]
onde $\Gamma\times S'$ denota a soma direta de $\Gamma$ cópias de $S'$, seja
um isomorfismo. (Supõe-se que em $\mathcal C$ existem as somas diretas que
intervêm aqui.) Seja $F$ um funtor contravariante de $\mathcal C$ para a
categoria dos grupos abelianos. Então $C^*(S'/S,F)$ é canonicamente
isomorfo ao grupo simplicial padrão de cocadeias homogêneas
$C^*(\Gamma,F(S'))$ e, portanto, $H^*(S'/S,F)$ é canonicamente isomorfo a
$H^*(\Gamma,F(S'))$.

\medskip\noindent
e. \result{EXEMPLO}{2} Seja $\mathcal C$ a categoria dos pré-esquemas. Denota-se
por $G_a$ («grupo aditivo») o funtor contravariante de $\mathcal C$ para a
categoria dos grupos abelianos definido por
\[
  G_a(X)=H^0(X,\mathcal O_X).
\]
Define-se de igual modo o funtor $G_m$ («grupo multiplicativo») por meio de
\[
  G_m(X)=H^0(X,\mathcal O_X)^*
\]
(grupo dos elementos invertíveis do anel $H^0(X,\mathcal O_X)$), e, mais
geralmente, o funtor $\operatorname{Gl}(n)$ («grupo linear de ordem $n$»)
mediante

\src{313}
\[
  \operatorname{Gl}(n)(X)=
  \operatorname{Gl}\bigl(n,H^0(X,\mathcal O_X)\bigr),
\]
que é um funtor de $\mathcal C$ à categoria dos grupos (não
necessariamente comutativos se $n>1$; para $n=1$ se recupera $G_m$). Além disso,
pode-se interpretar $\operatorname{Gl}(n)$ como um funtor de automorfismos
(cf. (c)) considerando a categoria fibrada $\mathcal F$ de base $\mathcal C$
tal que, para $X\in\mathcal C$, $\mathcal F_X$ seja a categoria dos feixes
localmente livres sobre $X$: com efeito, tem-se
\[
  \operatorname{Gl}(n)(X)=
  \operatorname{Aut}_{\mathcal F_X}(\mathcal O_X^n).
\]
De (b) segue-se que, se $\alpha:S'\longrightarrow S$ é um morfismo de
$\mathcal F$-descida (cf. parágrafo 2.c),
$H^1(S'/S,\operatorname{Gl}(n))$ contém o conjunto das classes (salvo
isomorfismo) de feixes localmente livres sobre $S$ cuja imagem inversa sobre
$S'$ é isomorfa a $\mathcal O_{S'}^n$, e essa inclusão é uma igualdade se e
somente se todo dado de descida sobre $\mathcal O_{S'}^n$ (relativo a
$\alpha:S'\longrightarrow S$) é efetivo. Quando $S$ é o espectro de um
anel local, isto significa, portanto,
$H^1(S'/S,\operatorname{Gl}(n))=(e)$, pois todo feixe localmente livre sobre $S$
é então trivial.

Assinalemos a equivalência das condições seguintes para um morfismo
$\alpha:S'\longrightarrow S$:
\begin{enumerate}
\item[(i)] o homomorfismo de aumento
\[
  H^0(S,\mathcal O_S)=G_a(S)\longrightarrow H^0(S'/S,G_a)
\]
é um isomorfismo;
\item[(ii)] $S'\longrightarrow S$ é um morfismo de $\mathcal F$-descida
(sendo $\mathcal F$ a categoria fibrada de base $\mathcal C$ considerada
acima).
\end{enumerate}
Se $S'\longrightarrow S$ é finito, essas condições equivalem também a:
\begin{enumerate}
\item[(iii)] $S'\longrightarrow S$ é um epimorfismo estrito
(cf. parágrafo 2.c).
\end{enumerate}

Suponhamos agora que $S=\operatorname{Spec}(A)$,
$S'=\operatorname{Spec}(A')$; então
\[
  C^n(S'/S,G_a)=C^n(A'/A,G_a)=\bigotimes_A^{n+1}A',
\]
sendo o operador coborde
$C^n(A'/A,G_a)\longrightarrow C^{n+1}(A'/A,G_a)$ a soma alternada dos
operadores de face
\[
  \partial_i(x_0\otimes x_1\otimes\cdots\otimes x_n)
  =x_0\otimes\cdots\otimes x_{i-1}\otimes1_{A'}\otimes
   x_i\otimes\cdots\otimes x_n.
\]
Do mesmo modo,
\[
  C^n(S'/S,G_m)=C^n(A'/A,G_m)
  \simeq\left(\bigotimes_A^{n+1}A'\right)^*,
\]
e as operações simpliciais em $C^*(A'/A,G_m)$ são induzidas pelas de
$C^*(S'/S,G_a)$. Explicitam-se analogamente as operações simpliciais em
$C^*(A'/A,\operatorname{Gl}(n))$. \emph{Em todos os casos que o
conferenciante conhece, tem-se}
\[
  H^i(A'/A,G_a)=0\qquad\text{para $i>0$},
\]
\emph{e, se $A$ é local, tem-se}
\[
  H^1(A'/A,G_m)=0
  \quad\text{e, mais geralmente,}\quad
  H^1(A'/A,\operatorname{Gl}(n))=(e)
\]
(quando $S'\longrightarrow S$ é um morfismo de $\mathcal F$-descida,
isto é, quando o diagrama
$A\longrightarrow A'\rightrightarrows A'\otimes_AA'$ é exato; compare-se com
o parágrafo 2.c). Observe-se que \emph{o «teorema 90» de Hilbert nada mais é
que a relação}
\src{314}
\[
  H^1(S'/S,G_m)=0
\]
\emph{quando $A$ é um corpo e $A'$ uma extensão de Galois finita deste
(cf. exemplo 1), e também se pode exprimir dizendo que, no caso
considerado, $S'\longrightarrow S$ é um morfismo de descida estrito para a
categoria fibrada dos feixes localmente livres de posto 1.} É sob esta
última forma que convém generalizar o teorema de Hilbert, variando tanto
as hipóteses sobre o morfismo $S'\longrightarrow S$ como as relativas aos
feixes quase-coerentes considerados.

Assinalemos, por fim, a equivalência das propriedades seguintes quando
$A$ é um anel local artiniano de ideal maximal $\mathfrak m$ e $A'$ uma
$A$-álgebra (denotando, para todo inteiro $k>0$, por $A_k$ (resp. $A'_k$) os
anéis $A/\mathfrak m^{k+1}$ (resp. $A'/\mathfrak m^{k+1}A'$)):
\begin{enumerate}
\item[(i)] $H^1(A'_k/A_k,G_a)=0$ para todo $k$.
\item[(ii)] $H^1(A'_k/A_k,G_m)=0$ para todo $k$.
\item[(iii)] $H^1(A'_k/A_k,\operatorname{Gl}(n))=(e)$ para todo $k$ e todo
$n$.
\end{enumerate}
Se $S'\longrightarrow S$ é um epimorfismo estrito, as condições
anteriores implicam inclusive que é um morfismo de descida estrito para os
módulos livres (de tipo finito ou não) sobre $A'$.

\result{OBSERVAÇÃO}{} A definição dos grupos $H^i(S'/S,G_m)$, quando
$S,S'$ são esquemas de corpos $A,A'$, se deve a Amitsur. O grupo
$H^2(S'/S,G_m)$ é particularmente interessante como variante «global» do
grupo de Brauer; para essa variante pode-se consultar [1], capítulo VII.

\fsection{B}{Descida por morfismos fielmente planos}

\subsection*{1. Enunciado dos teoremas de descida}
\addcontentsline{toc}{subsection}{1. Enunciado dos teoremas de descida}

\result{DEFINIÇÃO}{1.1} Um morfismo
$\alpha:S'\longrightarrow S$ de pré-esquemas diz-se \emph{plano} se, para todo
$x'\in S'$, $\mathcal O_{x'}$ é um módulo plano sobre o anel
$\mathcal O_{\alpha(x')}$ (isto é, se
$\mathcal O_{x'}\otimes_{\mathcal O_{\alpha(x')}}M$ é um funtor exato no
$\mathcal O_{\alpha(x')}$-módulo $M$). Um morfismo se diz
\emph{fielmente plano} se é plano e sobrejetivo.

Por exemplo, se $S=\operatorname{Spec}(A)$ e
$S'=\operatorname{Spec}(A')$, então $S'$ é plano sobre $S$ se e somente se
$A'$ é um $A$-módulo plano, e $S'$ é fielmente plano sobre $S$ se e somente se
$A'$ é um $A$-módulo fielmente plano (isto é, o funtor
$A'\otimes_A M$ no $A$-módulo $M$ é exato e fiel); isto significa também,
na terminologia de Serre [5], que o par $(A,A')$ é plano. Se $S'$ é
fielmente plano sobre $S$, então o funtor imagem inversa de feixes
quase-coerentes sobre $S$ é exato e fiel; em outras palavras, para que uma
sucessão de homomorfismos de feixes

\src{315}
quase-coerentes sobre $S$ seja exata, é necessário e suficiente que também o seja a sua
imagem inversa sobre $S'$ (em particular, para que um homomorfismo de feixes
quase-coerentes sobre $S$ seja um monomorfismo, resp. um epimorfismo, resp. um
isomorfismo, é necessário e suficiente que também o seja a sua imagem inversa sobre
$S'$). Essa propriedade continua válida se nos restringirmos a um aberto
qualquer de $S$ e, sob essa forma, caracteriza os morfismos fielmente
planos.

\result{DEFINIÇÃO}{1.2} Um morfismo $\alpha:S'\longrightarrow S$ diz-se
\emph{quase-compacto} se a imagem inversa de toda parte aberta quase-compacta
$U$ de $S$ é quase-compacta (isto é, união finita de abertos afins).

Evidentemente, basta verificar esta propriedade para os abertos afins de $S$.
Por exemplo, um morfismo afim (isto é, tal que a imagem inversa de um aberto
afim seja afim) é quase-compacto.

A classe dos morfismos planos, resp. fielmente planos, resp. quase-compactos,
é estável por composição e por «mudança de base», e contém, naturalmente,
os isomorfismos.

\result{TEOREMA}{1} Seja $\alpha:S'\longrightarrow S$ um morfismo de
pré-esquemas, fielmente plano e quase-compacto. Então $\alpha$ é um
\emph{morfismo de descida estrito} (A, definição 1.7) para a categoria
fibrada $\mathcal F$ dos feixes quase-coerentes (A, parágrafo 1, exemplo 2).

Esse enunciado significa duas coisas:
\begin{enumerate}
\item[(i)] Se $F$ e $G$ são dois feixes quase-coerentes sobre $S$, e $F'$ e
$G'$ suas imagens inversas sobre $S'$, então o homomorfismo natural
\[
  \operatorname{Hom}(F,G)\longrightarrow\operatorname{Hom}(F',G')
\]
é uma bijeção do primeiro membro sobre o subgrupo do segundo formado pelos
homomorfismos $F'\longrightarrow G'$ compatíveis com os dados de descida
canônicos sobre esses feixes, isto é, aqueles cujas imagens inversas pelas
duas projeções de $S''=S'\times_SS'$ sobre $S'$ dão um mesmo
homomorfismo $F''\longrightarrow G''$.
\item[(ii)] Todo feixe quase-coerente $F'$ sobre $S'$, provido de um dado de
descida relativo ao morfismo $\alpha:S'\longrightarrow S$
(A, definição 1.6), é isomorfo (provido deste dado) à imagem inversa de
um feixe quase-coerente $F$ sobre $S$.
\end{enumerate}
Tomando $F=\mathcal O_S$ em (i), obtém-se:

\result{COROLÁRIO}{1} Seja $G$ um feixe quase-coerente sobre $S$; sejam
$G'$ e $G''$ suas imagens inversas sobre $S'$ e sobre
$S''=S'\times_SS'$, e sejam $p_1,p_2$ as duas projeções de $S''$ sobre
$S'$. Então o seguinte diagrama de aplicações de conjuntos

\src{316}
\[
  \Gamma(G)\xrightarrow{\alpha^*}\Gamma(G')
  \underset{p_2^*}{\overset{p_1^*}{\rightrightarrows}}\Gamma(G'')
\]
é exato (A, definição 1.(a)).

Por outro lado, a conjunção de (i), (ii) e da definição 1.1 dá o

\result{COROLÁRIO}{2} Seja $G$ como no corolário 1. Então existe uma
correspondência biunívoca entre os subfeixes quase-coerentes de $G$ e os
subfeixes quase-coerentes de $G'$ cujas imagens inversas sobre $S''$, pelas
duas projeções $p_1,p_2$, dão o mesmo subfeixe de $G''$.

Naturalmente, tem-se um enunciado equivalente em termos de feixes quociente.
Como se viu (A, parágrafo 4, (e)), o teorema 1 deve ser considerado uma
generalização do «teorema 90» de Hilbert, e implica como casos particulares
diversas formulações em termos de 1-cohomologia. Para a demonstração, reduz-se
facilmente ao caso em que $S=\operatorname{Spec}(A)$,
$S'=\operatorname{Spec}(A')$, e para (i) reduz-se facilmente a demonstrar o
corolário 1, isto é, a exatidão do diagrama
\[
  M=A\otimes_A M\longrightarrow A'\otimes_A M
  \rightrightarrows A'\otimes_AA'\otimes_A M
\]
para todo $A$-módulo $M$, o que se deduz do lema mais geral:

\result{LEMA}{1.1} Seja $A'$ uma $A$-álgebra fielmente plana. Então, para
todo $A$-módulo $M$, o complexo aumentado por $M$
$C^*(A'/A,G_a)\otimes_A M$ (cf. A, parágrafo 4, (e)) é uma
\emph{resolução} de $M$.

Basta demonstrar que o complexo aumentado deduzido do anterior mediante a
mudança de base de $A$ para $A'$ satisfaz as mesmas conclusões. Isso conduz a
verificar o enunciado quando se substituem $A$ por $A'$ e $A'$ por
$A'\otimes_AA'$, o que nos reduz a um caso no qual existe um
homomorfismo de $A$-álgebras $A'\longrightarrow A$ (ou, em termos
geométricos, ao caso em que $S'$ sobre $S$ admite uma seção). Nesse caso,
isso se deduz das generalidades de A, parágrafo 4, (a). Assinalemos de passagem
o seguinte corolário, que generaliza um enunciado bem conhecido em
cohomologia de Galois (compare-se com A, parágrafo 4, (e)):

\result{COROLÁRIO}{} Se $A'$ é fielmente plana sobre $A$, tem-se
$H^0(A'/A,G_a)=A$ e
\[
  H^i(A'/A,G_a)=0\qquad\text{para $i\geq1$}.
\]

Para demonstrar a parte (ii) do teorema 1, procede-se como para (i),
reduzindo-se ao caso em que $S'$ sobre $S$ admite uma seção, onde se deduz
de (i) (cf. A, parágrafo 1, (c)).

Evidentemente, pode-se variar ad libitum o teorema 1 e seus corolários,
introduzindo diversas estruturas adicionais nos feixes (ou sistemas de
feixes) quase-coerentes considerados. Por exemplo, dar sobre $S$ um feixe
\src{317}
quase-coerente de álgebras comutativas «equivale» a dar sobre $S'$ um feixe desse
tipo provido de um dado de descida relativo a
$\alpha:S'\longrightarrow S$. Tendo em conta a correspondência funtorial
entre tais feixes quase-coerentes sobre $S$ e os pré-esquemas afins sobre
$S$, obtém-se a segunda afirmação do teorema seguinte:

\result{TEOREMA}{2} Seja $\alpha:S'\longrightarrow S$ como no teorema 1.
Então $\alpha$ é um \emph{morfismo de descida} (não estrito em geral)
(A, definição 2.4), e é um \emph{morfismo de descida estrito} para a
categoria fibrada dos esquemas afins sobre pré-esquemas
(A, definição 1.7).

A primeira afirmação do teorema significa o seguinte: sejam $X,Y$ dois
pré-esquemas sobre $S$; sejam $X',Y'$ suas imagens inversas sobre $S'$, e
$X'',Y''$ suas imagens inversas sobre $S''=S'\times_SS'$. Então o seguinte
diagrama de aplicações naturais
\[
  \operatorname{Hom}_S(X,Y)\xrightarrow{\alpha^*}
  \operatorname{Hom}_{S'}(X',Y')
  \underset{p_2^*}{\overset{p_1^*}{\rightrightarrows}}
  \operatorname{Hom}_{S''}(X'',Y'')
\]
é exato; isto é, $\alpha^*$ é uma bijeção de
$\operatorname{Hom}_S(X,Y)$ sobre a parte de
$\operatorname{Hom}_{S'}(X',Y')$ formada pelos homomorfismos compatíveis com
os dados de descida canônicos sobre $X',Y'$ (isto é, aqueles cujas
imagens inversas pelas duas projeções de $S''$ sobre $S'$ são iguais).
Isso se deduz facilmente do teorema 1, corolário 1, quando se restringe
a $Y$ afim sobre $S$; no caso geral, é preciso combinar o teorema 1 com o
resultado seguinte:

\result{LEMA}{1.2} Seja $\alpha:S'\longrightarrow S$ um morfismo fielmente plano
e quase-compacto. Então $S$ identifica-se com um \emph{espaço topológico
cociente} de $S'$; isto é, toda parte $U$ de $S$ tal que
$\alpha^{-1}(U)$ seja aberta é aberta.

Para completar o teorema 2, é preciso dar critérios de efetividade para um dado
de descida sobre um $S'$-pré-esquema $X'$ (quando não se supõe que $X'$ seja
afim sobre $S'$). Assinalemos primeiro que \emph{tal dado de descida não é
necessariamente efetivo}, mesmo se $S$ é o espectro de um corpo $k$, $S'$
o espectro de uma extensão quadrática $k'$ deste e $X'$ um esquema
algébrico próprio de dimensão 2 sobre $S'$ (como se pode ver, segundo Serre,
utilizando a superfície não projetiva de Nagata). \emph{Para que um dado de
descida sobre $X'/S'$ relativo a $\alpha:S'\longrightarrow S$ (fielmente
plano e quase-compacto) seja efetivo, é necessário e suficiente que $X'$ seja
união de abertos $X'_i$, afins sobre $S'$, que sejam «estáveis» pelo dado de
descida sobre $X'$.} Isso certamente se cumpre (quaisquer que sejam $X'/S'$
e o dado de descida sobre $X'$) se o morfismo
$\alpha:S'\longrightarrow S$ é \emph{radicial} (isto é, injetivo e com
extensões residuais que

\src{318}
são radiciais). Pode-se demonstrar também que continua a cumprir-se se
$\alpha:S'\longrightarrow S$ é \emph{finito} e toda parte finita de $X'$,
contida numa fibra de $X'$ sobre $S$, está contida num aberto de $X'$
afim sobre $S$ (esse é o \emph{critério de Weil}). Em particular, isso se
cumpre se $X'/S'$ é quase-projetivo, e nesse caso pode-se demonstrar que o
pré-esquema «descido» $X/S$ também é quase-projetivo (e projetivo se
$X'/S'$ também o é). Em resumo:

\result{TEOREMA}{3} Seja $\alpha:S'\longrightarrow S$ um morfismo de
pré-esquemas, fielmente plano e quase-compacto. Se $\alpha$ é \emph{radicial},
é um \emph{morfismo de descida estrito}. Se $\alpha$ é finito, é um
morfismo de descida estrito com respeito à categoria fibrada dos
pré-esquemas quase-projetivos (ou projetivos) sobre pré-esquemas.

\result{OBSERVAÇÕES}{} Ignoro se, na segunda afirmação anterior, a
hipótese de que $\alpha$ seja um morfismo \emph{finito} é realmente necessária;
em todo caso, verifica-se de maneira «formal» que pode ser substituída pela
seguinte hipótese, aparentemente mais fraca: \emph{para todo ponto de $S$
existe uma vizinhança aberta $U$, um $U'$ finito e fielmente plano sobre $U$, e um
$S$-morfismo de $U'$ em $S'$.} Um caso típico que não entra no anterior é
aquele em que $S=\operatorname{Spec}(A)$,
$S'=\operatorname{Spec}(\widehat A)$, onde $A$ é um anel local noetheriano
e $\widehat A$ seu completamento; ou também aquele em que $S'$ é quase-finito sobre
$S$ (isto é, localmente isomorfo a um aberto de um $S$-esquema finito) e não
finito. Nesses dois casos, o conferenciante também ignora a resposta à
pergunta seguinte: seja $X$ um $S$-esquema tal que $X'=X\times_SS'$ seja
projetivo sobre $S'$; é verdade que $X$ é projetivo sobre
$S$?\footnote{A errata de 1962 responde negativamente: um morfismo
$S'\longrightarrow S$ quase-finito, étale e sobrejetivo, ou um morfismo
$\operatorname{Spec}(\widehat A)\longrightarrow\operatorname{Spec}(A)$, não é
sempre um morfismo de descida estrito, mesmo se $A$ é o anel local de
uma curva algébrica sobre um corpo algebricamente fechado $k$ e
$S=\operatorname{Spec}(A)$. Pode-se encontrar assim um morfismo próprio e liso
$f:X\longrightarrow S$, que faz de $X$ um fibrado principal de base $S$ sob
uma curva elíptica $E$, tal que $f':X'\longrightarrow S'$ seja projetivo, mas
$f$ não o seja. Portanto, também é um exemplo de fibrado principal homogêneo
não isotrivial sob um esquema abeliano.}

\subsection*{2. Aplicação à descida de certas propriedades dos morfismos}
\addcontentsline{toc}{subsection}{2. Descida de propriedades dos morfismos}

Seja $P$ uma classe de morfismos de pré-esquemas. Seja
$\alpha:S'\longrightarrow S$ um morfismo de pré-esquemas, seja
$f:X\longrightarrow Y$ um morfismo de $S$-pré-esquemas e seja
$f':X'\longrightarrow Y'$ a imagem inversa de $f$ por $\alpha$. Cabe
perguntar-se então se a relação «$f'\in P$» implica «$f\in P$». A resposta é
afirmativa em muitos casos importantes quando se supõe que $\alpha$ é
\emph{fielmente plano e quase-compacto} (sendo às vezes supérflua essa última
hipótese). Isso se vê diretamente e sem dificuldade se $P$ é a classe dos
morfismos sobrejetivos, resp. radiciais (esses casos se deduzem da
sobrejetividade de $\alpha$), resp. planos, resp. fielmente planos, resp. lisos
(esses casos se deduzem da planicidade fiel de $\alpha$), resp. de tipo finito.
Utilizando os teoremas 1 e 2 e o lema 1.2, vê-se também que ocorre o mesmo se
$P$ é uma das classes seguintes: isomorfismos, imersões abertas, imersões
fechadas, imersões (se $f$ é de tipo finito e $Y$
localmente noetheriano), morfismos afins, morfismos finitos, morfismos
quase-finitos, morfismos abertos, morfismos fechados, homeomorfismos, morfismos
separados,

\src{319}
morfismos próprios. O único caso importante não elucidado é o dos
morfismos projetivos e quase-projetivos, já assinalado na observação do
parágrafo 1.

\subsection*{3. Descida por morfismos finitos fielmente planos}
\addcontentsline{toc}{subsection}{3. Descida por morfismos finitos fielmente planos}

Seja $\alpha:S'\longrightarrow S$ um morfismo finito correspondente a um feixe de
álgebras $\mathcal A'$ sobre $S$ que, como feixe de módulos, é localmente livre
de tipo finito e é $\ne0$ em toda parte; então $\alpha$ é um morfismo
fielmente plano e quase-compacto, ao qual se podem aplicar os resultados
anteriores. Dar um feixe quase-coerente $F'$ sobre $S'$ equivale a dar o feixe
quase-coerente $\alpha_*(F')$ sobre $S$, provido de sua estrutura de
$\mathcal A'$-módulo (observando que
$\mathcal A'=\alpha_*(\mathcal O_{S'})$). Para simplificar, este feixe sobre $S$
também será denotado por $F'$. As duas imagens inversas $p_i^*(F')$ de $F'$
sobre $S'\times_SS'$ correspondem igualmente aos feixes quase-coerentes de
$(\mathcal A'\otimes_{\mathcal O_S}\mathcal A')$-módulos
\[
  F'\otimes_{\mathcal O_S}\mathcal A'
  \quad\text{e}\quad
  \mathcal A'\otimes_{\mathcal O_S}F'.
\]
Dar um $(S'\times_SS')$-homomorfismo do primeiro ao segundo equivale a dar um
homomorfismo de $(\mathcal A'\otimes_{\mathcal O_S}\mathcal A')$-módulos e,
tendo em conta que $\mathcal A'$ é localmente livre, isto equivale a dar
um homomorfismo de $(\mathcal A'\otimes\mathcal A')$-módulos:
\[
  \mathcal U=
  \operatorname{Hom}_{\mathcal O_S}(\mathcal A',\mathcal A')
  =\mathcal A'\otimes\check{\mathcal A}'
  \longrightarrow
  \operatorname{Hom}_{\mathcal O_S}(F',F'),
\]
isto é, a dar, para cada seção $\xi$ de $\mathcal U$ sobre um aberto $V$,
um homomorfismo de $\mathcal O_S$-módulos
$T_\xi:F'|_V\longrightarrow F'|_V$ que satisfaz as condições
\begin{equation}
  T_{f\xi}(x)=fT_\xi(x),
  \qquad T_{\xi f}(x)=T_\xi(fx),
  \tag{3.1}
\end{equation}
onde $f,x$ são, respectivamente, seções de $\mathcal A',F'$ sobre um
aberto de $S$ contido em $V$. As condições (i) e (ii) de um dado de
descida (A, parágrafo 1, (c)) escrevem-se então, respectivamente:
\begin{equation}
  T_{1_{\mathcal U}}x=x,
  \qquad\text{isto é,}\qquad
  T_{1_{\mathcal U}}=\operatorname{id}_{F'},
  \tag{3.2}
\end{equation}
\begin{equation}
  T_{\xi\eta}=T_\xi T_\eta.
  \tag{3.3}
\end{equation}
Em outras palavras, \emph{um dado de descida sobre $F'$ equivale a uma
representação do feixe de $\mathcal O_S$-álgebras
$\mathcal U=\operatorname{Hom}_{\mathcal O_S}(\mathcal A',\mathcal A')$ no
feixe de $\mathcal O_S$-álgebras
$\operatorname{Hom}_{\mathcal O_S}(F',F')$ que satisfaz as duas condições de
linearidade (3.1).} Se se tem um emparelhamento de feixes quase-coerentes sobre
$S'$:
\[
  F'_1\times F'_2\longrightarrow F'_3
\]
(que pode ser interpretado como um emparelhamento de feixes de
$\mathcal A'$-módulos sobre $S$),

\src{320}
e dados de colagem sobre os $F'_i$, definidos por homomorfismos
\[
  T_i:\mathcal U\longrightarrow
  \operatorname{Hom}_{\mathcal O_S}(F'_i,F'_i)
  \qquad (i=1,2,3),
\]
então esses dados são compatíveis com o emparelhamento dado, no sentido
evidente do termo, se e somente se vigora a condição seguinte.

Para toda seção $\xi$ de $\mathcal U$ sobre um aberto, denotando por
\[
  \Delta\xi=\sum_i\xi'_i\otimes_{\mathcal A'}\xi''_i
\]
a seção de $\mathcal U\otimes_{\mathcal A'}\mathcal U$ (considerando
$\mathcal U$ como $\mathcal A'$-módulo mediante sua estrutura esquerda)
definida pela fórmula
\[
  \xi\mathbin{\cdot}(fg)=\sum_i\xi'_i(f)\xi''_i(g)
\]
(onde $f$ e $g$ são duas seções de $\mathcal A'$ sobre um aberto menor),
tem-se a fórmula
\begin{equation}
  T^{(3)}_\xi(x\mathbin{\cdot}y)
  =\sum_iT^{(1)}_{\xi'_i}x\mathbin{\cdot}T^{(2)}_{\xi''_i}y
  \tag{3.4}
\end{equation}
para todo par de seções $x,y$ de $F'_1,F'_2$, respectivamente, sobre um
aberto menor.\footnote{O texto impresso diz «de $\mathcal A'$». Os domínios de
$T^{(1)}$ e $T^{(2)}$, assim como o emparelhamento
$F'_1\times F'_2\longrightarrow F'_3$, obrigam a que $x$ e $y$ sejam seções
de $F'_1$ e $F'_2$, respectivamente.} (Essa propriedade pode ser expressa
dizendo que os homomorfismos $T^{(i)}$ são compatíveis com a aplicação
diagonal de $\mathcal U$ com respeito ao emparelhamento dado.) Em particular, as
fórmulas (3.1) a (3.4) permitem interpretar, em termos de representações de
álgebras com aplicações diagonais, os dados de descida sobre um feixe
quase-coerente de álgebras sobre $S'$ e, portanto, também (restringindo-se às
álgebras comutativas) os dados de descida sobre um $S'$-esquema afim.

Daí se passa a uma interpretação análoga dos dados de descida sobre um
$S'$-pré-esquema $X'$ qualquer: dar tal $X'$ equivale a dar um pré-esquema
$X'$ sobre $S$, provido de um homomorfismo de $\mathcal O_S$-álgebras
\[
  \mathcal A'\longrightarrow\mathcal O_{X'},
\]
e um dado de descida sobre $X'$ equivale a dar um homomorfismo de feixes
(compatível com o morfismo $h:X'\longrightarrow S'$)
\[
  \mathcal U\longrightarrow
  \operatorname{Hom}_{h^{-1}(\mathcal O_S)}
  (\mathcal O_{X'},\mathcal O_{X'}),
\]
que satisfaz condições análogas às condições (3.1) a (3.4) anteriores.

\textsc{Exemplo 1 («Weil»).} -- Suponhamos que $S'/S$ é um recobrimento étale
de Galois de grupo de Galois $\Gamma$ (cf. A, parágrafos 3 e 4.d). Então, um
dado de descida sobre um feixe quase-coerente $F'$ sobre $S'$ (resp. sobre um
$S'$-pré-esquema $X'$) equivale a dar uma representação de $\Gamma$ por
automorfismos de $(S',F')$ (resp. de $(S',X')$) compatível com as operações
de $\Gamma$ sobre $S'$. Este resultado

\src{321}
é «formal», isto é, demonstra-se em termos de categorias, mas, do
ponto de vista deste número, deduz-se também da estrutura explícita de
$\mathcal U$ (provido de sua estrutura de anel, do homomorfismo de anéis
$\mathcal A'\longrightarrow\mathcal U$ e da aplicação diagonal), conhecida
por completo graças ao resultado seguinte: $\mathcal U$ admite, como
$\mathcal A'$-módulo esquerdo, uma base formada pelas seções de
$\mathcal U$ que correspondem aos elementos de $\Gamma$.

\textsc{Exemplo 2 («Cartier»).} -- Seja $p$ um número primo; suponhamos que
$p\mathcal O_S=0$ (isto é, $\mathcal O_S$ é de característica $p$),
$(\mathcal A')^p\subset\mathcal O_S=\mathcal A$ (isto é, $S'/S$ é radicial
de altura 1) e que o feixe de álgebras $\mathcal A'$ sobre $\mathcal A$ admite
localmente uma $p$-base, isto é, uma família $(x_i)$ de seções tal que
$\mathcal A'$ está gerada como álgebra pelos $x_i$, submetidos às únicas
condições $x_i^p=0$. Supomos que o conjunto dos $i$ é finito, de
cardinal $n$. Seja $\mathfrak g$ o feixe das $\mathcal A$-derivações de
$\mathcal A'$; é um feixe de $\mathcal A'$-módulos localmente livre de posto
$n$ e, além disso, um feixe de $p$-álgebras de Lie sobre $\mathcal A$ (mas não sobre
$\mathcal A'$) que satisfaz a condição
\begin{equation}
  [X,fY]=X(f)Y+f[X,Y].
  \tag{3.5}
\end{equation}

\textsc{Lema.} -- $\mathcal U=
\operatorname{Hom}_{\mathcal O_S}(\mathcal A',\mathcal A')$ está gerada,
como $\mathcal O_S$-álgebra provida de um homomorfismo de álgebras
$\mathcal A'\longrightarrow\mathcal U$, pelo sub-$\mathcal A'$-módulo
esquerdo $\mathfrak g$, com as relações adicionais
\begin{equation}
  \left\{
  \begin{aligned}
    Xf-fX&=X(f),\\
    XY-YX&=[X,Y],\\
    X^p&=X^{(p)}.
  \end{aligned}
  \right.
  \tag{3.6}
\end{equation}

Do lema anterior se deduz que um dado de descida sobre o feixe
quase-coerente $F'$ sobre $S'$ equivale a dar, para todo $X\in\mathfrak g$, um
$\mathcal O_S$-endomorfismo $\overline X$ de $F'$ que satisfaz as condições
\begin{align}
  \overline{fX}&=f\overline X, \tag{3.7}\\
  \overline X(fx)&=X(f)x+f\overline X(x), \tag{3.8}\\
  \overline{[X,Y]}&=[\overline X,\overline Y], \tag{3.9}\\
  \overline{X^{(p)}}&=\overline X^p. \tag{3.10}
\end{align}
(Isso é o que se poderia chamar uma conexão linear sobre $F'$, sem curvatura
e compatível com a potência $p$-ésima.) Explicita-se de igual modo a noção
de dado de descida sobre um $S'$-pré-esquema $X'$; a relação (3.4) é
substituída aqui pela condição de que os $\overline X$ sejam derivações de
$\mathcal O_{X'}$. Como o morfismo $S'\longrightarrow S$ é radicial, o
teorema 3 garante que todo dado de descida desse tipo é

\src{322}
efetivo e, portanto, define um $S$-pré-esquema $X$.

Observe-se que não foi necessário impor nenhuma hipótese de planicidade, de
não singularidade nem de finitude de tipo algum sobre $F'$, resp. $X'$.

\subsection*{4. Aplicação a critérios de racionalidade}
\addcontentsline{toc}{subsection}{4. Aplicação a critérios de racionalidade}

Seja $X$ um $S$-pré-esquema tal que a imagem direta de $\mathcal O_X$ sobre
$S$ seja $\mathcal O_S$; essa propriedade continuará válida após toda mudança
plana $S'\longrightarrow S$ da base $S$. Se $F$ é um feixe invertível
(isto é, localmente livre de posto 1) sobre $X$, os automorfismos de $F$, que se
identificam com as seções invertíveis de $\mathcal O_X$, correspondem
biunívocamente às seções invertíveis de $\mathcal O_S$. Seja então $s$
uma seção de $X$ sobre $S$; chamaremos, de maneira gráfica, seção de $F$
sobre $s$ a uma seção do feixe invertível $s^*(F)$ sobre $S$. Do anterior
deduz-se que, se $F_i$ ($i=1,2$) são dois feixes invertíveis sobre $X$, cada um
provido de uma seção sobre $s$, e se $F_1$ e $F_2$ são isomorfos, existe um
único isomorfismo de $F_1$ sobre $F_2$ compatível com as seções em questão
(isto é, que transforma a primeira na segunda). Por outro lado, e
independentemente da seção $s$, convenhamos em considerar equivalentes
dois feixes invertíveis $F_1$ e $F_2$ sobre $X$ tais que todo ponto de $S$ tenha
uma vizinhança aberta $U$ para a qual as restrições de $F_1$ e $F_2$ a
$X|U$ sejam isomorfas. Então todo feixe invertível $F$ sobre $X$ é equivalente
a um feixe invertível $F_1$ provido de uma seção distinguida sobre $s$
(toma-se $F_1=F\,s^*(F)^{-1}$), e $F_1$ está determinado salvo isomorfismo. Em
outras palavras, a classificação dos feixes invertíveis sobre $X$ salvo
equivalência é a mesma que a classificação, salvo isomorfismo, dos feixes
invertíveis providos de uma seção distinguida.

Como essas propriedades continuam válidas após uma mudança plana
$\alpha:S'\longrightarrow S$ da base (substituindo a seção $s$ por sua
imagem inversa $s'$ por $\alpha$), se conclui, tendo em conta o teorema 1:

\emph{Seja o pré-esquema $X/S$ como acima e suponhamos que admita uma seção
$s$; seja $\alpha:S'\longrightarrow S$ um morfismo fielmente plano e
quase-compacto; seja $F'$ um feixe invertível sobre $X'=X\times_SS'$. Para que $F'$
seja equivalente à imagem inversa sobre $X'$ de um feixe invertível $F$ sobre
$X$, é necessário e suficiente que suas imagens inversas $p_1^*(F')$ e
$p_2^*(F')$ sobre $X'\times_XX'=X\times_S(S'\times_SS')$ sejam equivalentes. Se
for assim, $F$ está determinado salvo equivalência.} (Dir-se-á então que $F'$
é racional sobre $S$.)

Inspirando-se nesse princípio quando $\alpha:S'\longrightarrow S$ é como no
exemplo 1 ou no exemplo 2 do número anterior, obtêm-se os critérios de
racionalidade de Weil ou de Cartier. (Observe-se que esses autores se limitam ao
caso em que $S$ e $S'$

\src{323}
são espectros de corpos; a fortiori, $S$ é então o espectro de um anel
local, e a relação de equivalência introduzida acima nada mais é que a relação
de isomorfia.) No primeiro caso, $F'$ é racional sobre $S$ se e somente se seus
transformados por $\Gamma$ são equivalentes a $F'$. Para exprimir o critério
de racionalidade no segundo caso, consideram-se, em geral, o morfismo
diagonal
\[
  X'\longrightarrow X''=X'\times_XX'
\]
de $X'/X$, o feixe de ideais $\mathcal I$ correspondente sobre
$X'\times_XX'$ e o feixe $\mathcal I/\mathcal I^2$, que se identifica com sua
imagem inversa $\Omega^1_{X'/X}$ sobre $X'$. Como as restrições dos
$F''_i=p_i^*(F')$ ($i=1,2$) à diagonal são isomorfas (pois são isomorfas a
$F'$), isto é, $F''_1(F''_2)^{-1}=F''$ tem uma restrição à diagonal
que é trivial, segue-se que a restrição de $F''$ a
$(X'',\mathcal O_{X''}/\mathcal I^2)$ é dada, salvo isomorfismo, por um
elemento bem determinado $\xi$ de
\[
  H^1(X'',\mathcal I/\mathcal I^2)
  =H^1(X',\Omega^1_{X'/X}).
\]
Além disso, nesse caso tem-se
\[
  \Omega^1_{X'/X}
  =\Omega^1_{S'/S}\otimes_{\mathcal O_S}\mathcal O_X,
\]
e, portanto, se $\Omega^1_{S'/S}$ é localmente livre sobre $S$ (como no
caso de Cartier), $\xi$ define uma seção de
\[
  R^1f'_*(\mathcal O_{X'})\otimes\Omega^1_{S'/S}
\]
sobre $S'$ (chamada classe de Atiyah--Cartier do feixe invertível $F'$ sobre
$X'/S$), cuja anulação é necessária e suficiente para que sejam equivalentes
as imagens inversas de $F'$ pelas duas projeções de
\[
  (X'',\mathcal O_{X''}/\mathcal I^2)
  =X\times_S(S'',\mathcal O_{S''}/\mathcal J^2)
\]
sobre $X'$ (onde $\mathcal J$ é o feixe de ideais sobre
$S''=S'\times_SS'$ definido pelo morfismo diagonal
$S'\longrightarrow S'\times_SS'$). Essa anulação é, pois, trivialmente
necessária para que sejam equivalentes as imagens inversas de $F'$ sobre o
próprio $X''=X\times_SS''$ e, portanto, também para que $F'$ seja equivalente
à imagem inversa de um feixe invertível $F$ sobre $X$. Além disso, a classe de
Atiyah--Cartier pode ser interpretada também como a obstrução à existência,
localmente sobre $S'$, de uma conexão de $F'$ relativa às derivações de
$X'/X$; tal conexão fica, além disso, determinada quando se conhecem as
derivações de $F'$ correspondentes às extensões naturais a $X'$ das
derivações de $S'/S$. Disso e dos desenvolvimentos do número anterior
conclui-se facilmente que, no caso do exemplo 2 desse número e quando $X/S$
admite uma seção, a anulação da classe de Atiyah--Cartier é também
suficiente para que $F'$ seja racional sobre $S$.

\subsection*{5. Aplicação à restrição da base em um esquema abeliano}
\addcontentsline{toc}{subsection}{5. Restrição da base em um esquema abeliano}

Seja $S$ um pré-esquema. Chama-se esquema abeliano sobre $S$ um esquema $X$
liso e próprio

\src{324}
sobre $S$ cujas fibras nos pontos $x\in S$ são variedades abelianas sobre os
$\kappa(x)$. Suponhamos que $S$ seja noetheriano e regular
(isto é, que seus anéis locais sejam regulares); então, utilizando o
teorema de conexidade de Murre [4] (ao menos no caso de «características
iguais», no qual o teorema citado está atualmente demonstrado), pode-se
demonstrar que toda seção racional de $X$ sobre $S$ está definida em toda
parte (isto é, é uma seção), o que generaliza um teorema clássico de
Weil. Deduz-se, mais geralmente, que, se $X'$ é um esquema liso sobre $S$,
então toda $S$-aplicação racional de $X'$ em $X$ está definida em toda
parte. Disso se obtém o seguinte, que generaliza um resultado de
Chow--Iang:\footnote{O texto impresso diz «Igusa--Lang»; a errata de 1962
prescreve «Chow--Iang».} Seja $S$ noetheriano e regular, e denote $K$ seu
anel de funções racionais (produto direto de corpos); seja $X$ um
esquema abeliano sobre $K$. Se $X$ é isomorfo a um $K$-esquema da forma
$X_0\times_S\operatorname{Spec}(K)$, onde $X_0$ é um esquema abeliano sobre
$S$, então $X_0$ está determinado salvo isomorfismo único.

Utilizando o resultado de unicidade anterior, vê-se que a questão da
restrição da base em $X$ é local sobre $S$ (e, portanto, que basta saber
efetuar a restrição aos $\operatorname{Spec}(\mathcal O_x)$, com
$x\in S$). Vê-se de igual maneira que, se
$\alpha:S'\longrightarrow S$ é um morfismo liso de tipo finito, se $K'$ é o
anel de funções racionais de $S'$, e se $X\otimes_KK'$ é da forma
$X'_0\times_{S'}\operatorname{Spec}(K')$, então $X'_0$ está provido de um
dado de descida canônico relativo a $\alpha$. Tendo em conta o teorema 3,
conclui-se:

\result{PROPOSIÇÃO}{5.1} Seja $S$ um pré-esquema noetheriano, regular e
irredutível, de corpo de funções racionais $K$; seja $K'$ uma extensão
finita de $K$, não ramificada sobre $S$; seja $S'$ a normalização de $S$ em
$K'$ (que é, portanto, um recobrimento étale de $S$); e seja $X$ um esquema
abeliano sobre $K$ tal que $X\otimes_KK'$ seja da forma
$X'_0\times_{S'}\operatorname{Spec}(K')$, onde $X'_0$ é um esquema abeliano
projetivo sobre $S'$. Então $X$ é da forma
$X_0\times_S\operatorname{Spec}(K)$, onde $X_0$ é um esquema abeliano
projetivo sobre $S$.

\result{OBSERVAÇÕES}{} O conferenciante ignora se pode ser substituída a
hipótese de que $S'\longrightarrow S$ seja um recobrimento étale sobrejetivo
(que permite utilizar o teorema 3) pela hipótese de que seja um morfismo
liso, sobrejetivo e de tipo finito (inclusive se se supõe que é um morfismo
étale), ou se a proposição continua válida sem supor que $X'_0$ seja
projetivo sobre $S'$ (condição que talvez se cumpra automaticamente).

\subsection*{6. Aplicação a critérios de trivialidade local e isotrivialidade}
\addcontentsline{toc}{subsection}{6. Critérios de trivialidade local e isotrivialidade}

Sejam $S$ um pré-esquema, $G$ um «pré-esquema em grupos» sobre $S$, e $P$ um
pré-esquema sobre $S$ sobre o qual «atua $G$» (pela direita). Diz-se que $P$
é formalmente principal homogêneo sob $G$ se o morfismo bem conhecido
\[
  G\times_SP\longrightarrow P\times_SP
\]

\src{325}
induzido pela ação de $G$ sobre $P$ é um isomorfismo. Doravante, suporemos que
$G$ seja plano sobre $S$ (e, portanto, fielmente plano sobre $S$), e reservaremos
o nome de fibrado principal homogêneo sob $G$ para um fibrado formalmente
principal homogêneo $P$ que seja fielmente plano e quase-compacto sobre $S$.
É imediato que isso equivale a dizer que se pode encontrar uma mudança
fielmente plana e quase-compacta
$S'\longrightarrow S$ da base $S$, tal que o fibrado formalmente principal
homogêneo
\[
  P'=P\times_SS'
\]
sob $G'=G\times_SS'$ seja trivial, isto é, isomorfo a $G'$ (isto é, que
admita uma seção); em particular, pode-se tomar $S'=P$. Observe-se, além disso,
que, se $S$ é localmente noetheriano, a hipótese de planicidade fiel sobre $P$
equivale à hipótese de que
\[
  \widehat P_s=P\times_S\operatorname{Spec}(\widehat{\mathcal O}_s)
\]
seja fielmente plano sobre $\widehat{\mathcal O}_s$ para todo $s\in S$ (onde
$\widehat{\mathcal O}_s$ denota o completamento do anel local
$\mathcal O_s$), como se deduz do fato de que $\widehat{\mathcal O}_s$ é
fielmente plano sobre $\mathcal O_s$. Além disso, se $P$ é de tipo finito sobre
$S$ localmente noetheriano, o conjunto dos pontos $s$ que satisfazem a
condição anterior é construtível; portanto, se $S$ é um «pré-esquema de
Jacobson» (por exemplo, um esquema de tipo finito sobre um corpo ou, mais
geralmente, o espectro de um anel de Jacobson), basta verificar a
condição em questão nos pontos fechados de $S$. Isso nos reduz ao caso em
que a base é o espectro de um anel local completo $A$. Quando
$S=\operatorname{Spec}(A)$ ($A$ um anel local noetheriano completo) e $P$ é
de tipo finito sobre $S$, a planicidade fiel de $P/S$ equivale também à
existência de um $S'$ finito e plano sobre $S$ tal que $P'$ seja trivial; e, se
além disso $G$ é liso sobre $S$, pode-se supor que $S'$ seja étale sobre $S$.
Por conseguinte, se além disso o corpo residual de $A$ é algebricamente fechado
(«caso geométrico»), $P$ é fielmente plano sobre $A$ se e somente se é trivial.
Assim, pois, se $S$ é um pré-esquema algébrico sobre um corpo algebricamente
fechado e $G$ é liso e de tipo finito sobre $S$, vê-se que a condição de
planicidade fiel sobre $S$ equivale à condição de trivialidade analítica (SLF)
de Serre ([6], pp. 1--12).

Podem-se introduzir outros tipos, mais fortes, de condições sobre $P$ que
tenham a natureza de uma «trivialidade local». Diremos, em particular, que
$P$ é isotrivial (resp. estritamente isotrivial) se, para todo $s\in S$,
existem uma vizinhança aberta $U$ de $s$ e um morfismo finito e fielmente plano
(resp. um recobrimento étale sobrejetivo) $U'\longrightarrow U$ tais que
$P'=P\times_SU'$ seja trivial. (Nossa terminologia difere da de Serre [6], que
chama localmente isotrivial ao que nós chamamos estritamente isotrivial.) A
isotrivialidade estrita é especialmente útil se $G$ é liso sobre $S$, mas, em
contrapartida, é uma noção inadequada nos demais casos.

Se $G$ é afim sobre $S$, todo fibrado principal homogêneo $P$ sob $G$ é afim
pelo parágrafo 2, de onde, graças ao teorema 2, resulta a possibilidade de
«descer»

\src{326}
tais fibrados por morfismos fielmente planos e quase-compactos. Tomemos, em
particular, $G=\operatorname{Gl}(n)_S$, definido pela condição de que o
funtor de $S$-pré-esquemas
\[
  S'\longmapsto\operatorname{Hom}_S(S',G)
\]
(com valores na categoria dos grupos) se identifique com o funtor
\[
  \operatorname{Gl}(n)(S')
  =\operatorname{Gl}\bigl(n,H^0(S',\mathcal O_{S'})\bigr)
\]
de A, parágrafo 4.e. Utilizando os fatos (i) de que todo fibrado principal
homogêneo sob $G$ (resp. todo feixe localmente livre de posto $n$ sobre $S$) se
torna isomorfo ao objeto «trivial» $G$ (resp. $\mathcal O_S^n$) após uma mudança
fielmente plana e quase-compacta adequada de $S$; (ii) de que podem ser descidos
os objetos do tipo considerado (fibrados principais homogêneos sob $G$,
resp. feixes localmente livres de posto $n$) por tais morfismos; e, por fim,
(iii) de que o grupo de automorfismos do fibrado trivial sobre um $S'/S$ é
funtorialmente isomorfo ao grupo de automorfismos do feixe localmente livre
trivial de posto $n$ sobre $S'$, conclui-se «formalmente» que «dá no mesmo»
dar sobre $S$ (ou sobre um $S'/S$) um fibrado principal homogêneo de grupo $G$
ou dar aí um feixe localmente livre de posto $n$. (Mais precisamente, tem-se
uma equivalência de categorias fibradas.) Conclui-se, em particular:

\result{PROPOSIÇÃO}{6.1} Todo fibrado principal homogêneo de grupo
$\operatorname{Gl}(n)_S$ é localmente trivial.

Mediante argumentos conhecidos obtém-se o mesmo resultado para grupos
estruturais tais como $\operatorname{Sl}(n)_S$, $\operatorname{Sp}(n)_S$ e
produtos de tais grupos. Deduz-se também que, se $F$ é um subgrupo fechado
de $G=\operatorname{Gl}(n)_S$, plano sobre $S$, tal que existe o quociente
$G/F$, e se $G$ é um fibrado principal homogêneo isotrivial (resp.
estritamente isotrivial) sobre $G/F$, de grupo $F\times_S(G/F)$, então todo
fibrado principal homogêneo de grupo $F$ é isotrivial (resp. estritamente
isotrivial). Isso se aplica a todos os «grupos lineares» sobre $S$ utilizados
até agora e, em particular, quando $G=S\times_k\Gamma$, onde $S$ é um
pré-esquema sobre o corpo $k$ e $\Gamma$ um grupo linear no sentido clássico
e, em particular, liso, sobre $k$. Isso resolve, pois, para tais grupos, uma
questão de Serre (loc. cit.).

Assinalemos igualmente que, para a maioria dos grupos lisos sobre $S$ conhecidos,
sejam lineares ou não, e em todo caso para os da forma
$S\times_k\Gamma$ anterior, pode demonstrar-se que todo fibrado principal
homogêneo isotrivial é estritamente isotrivial. Isso resolve, em particular,
outra questão de Serre (loc. cit., pp. 1--14), dado que um fibrado principal
homogêneo obtido por uma descida à Cartier (cf. parágrafo 3, exemplo 2) é
manifestamente isotrivial.

\result{OBSERVAÇÃO}{} Uma das dificuldades essenciais nessas questões
(além da questão da existência dos esquemas quociente $G/F$) é a
falta de critérios de efetividade para um dado de descida por um morfismo
fielmente plano não finito.

\src{327}
\begin{center}
  \textbf{BIBLIOGRAFIA}
\end{center}

\begin{description}
  \item[{[1]}] Dieudonné (J.) e Grothendieck (Alexander). --- \emph{Éléments
    de géométrie algébrique}, a aparecer nas Publications
    mathématiques de l'Institut des Hautes Études Scientifiques.
  \item[{[2]}] Grauert (H.) und Remmert (R.). --- \emph{Komplexe Räume}, Math.
    Annalen, t. 136, 1958, pp. 245--318.
  \item[{[3]}] Grothendieck (Alexander). --- \emph{Géométrie formelle et
    géométrie algébrique}. Séminaire Bourbaki, t. 11, 1958/59, n.º 182.
  \item[{[4]}] Murre (J. P.). --- \emph{On a connectedness theorem for a
    birational transformation at a simple point}, Amer. J. Math., t. 80, 1958,
    pp. 3--15.
  \item[{[5]}] Serre (Jean-Pierre). --- \emph{Géométrie algébrique et géométrie
    analytique}, Ann. Inst. Fourier Grenoble, t. 6, 1955--56, pp. 1--42.
  \item[{[6]}] Serre (Jean-Pierre). --- \emph{Espaces fibrés algébriques},
    Séminaire Chevalley, t. 2, 1958: Anneaux de Chow et applications, n.º 1.
\end{description}
